\documentclass[12pt]{article}
\usepackage[utf8]{inputenc}
\usepackage[T1]{fontenc}
\usepackage{lmodern}
\usepackage[margin=1.5in]{geometry}
\usepackage{setspace}
\usepackage{titling}

\onehalfspacing

\title{\textbf{The Thorns}}
\author{NASA Mission Series v1.23 --- Mercury-Atlas 7 / Aurora 7\\
\small A story about a plant that charges admission and a pilot who paid it}
\date{}

\begin{document}
\maketitle

The devil's club grows in the wet ravines where the light never quite reaches the ground. You find it where the stream cuts through old growth, where the Douglas-fir canopy is so thick that noon looks like dusk, where the soil is black with centuries of decomposed needles and the air smells like earth and rain and something green and ancient.

You see the leaves first. They are enormous --- broader than your hand, broader than a dinner plate, shaped like a maple leaf that decided to be ambitious. They catch what light filters through the canopy the way a satellite dish catches signal: broad, angled, optimized for collection in conditions of scarcity.

Then you see the spines.

\begin{center}$\ast\quad\ast\quad\ast$\end{center}

On May 24, 1962, a man named Scott Carpenter orbited the Earth three times in a capsule called Aurora 7. He tapped the hull. A shower of particles erupted from the spacecraft's surface. Ice crystals, sublimating at orbital dawn. The firefly mystery, solved.

The solving cost fuel. By the end of orbit two, 75\% of the manual attitude control fuel was consumed. The retrofire was three seconds late. The spacecraft was yawed twenty-five degrees. $\cos(25°) = 0.906$. Ninety-one percent of the thrust was correct. Nine percent was sideways. And nine percent became four hundred and two kilometers of overshoot.

\begin{center}$\ast\quad\ast\quad\ast$\end{center}

The devil's club does not apologize for its thorns. The thorns are not a flaw in the plant's design. They are the design. The armed all-healer: \textit{Oplopanax horridus}. The horribly bristly armed all-healer.

The thorns and the medicine grow from the same stem. The curiosity and the overshoot share the same fuel tank. You choose to reach in, or you don't. But you don't get to choose which parts of the plant you touch.

\vspace{2em}
\begin{flushright}
\textit{for Booker T. Washington, born April 5, 1856}\\
\textit{who built an institution from nothing}
\end{flushright}

\end{document}
